\exercisesfor{5}

在下列习题中，假设 §5 的假设和记号成立。F 表示自由群 \(\mathrm{F}(\mathbf{X})\)，g 表示从 F 到 Magnus 群 \(\Gamma(\mathbf{X})\) 的唯一同态，满足 \(g(x)=1+x\) 对所有 \(x\in X\) 成立（参见~定理 1）。

\begin{enumerate}
\def\labelenumi{\arabic{enumi}.}
\tightlist
\item
  给代数 \(\mathbf{K}[\mathbf{F}]\)（\(\mathbf{F}\) 的代数）赋予由 \((\mathbf{I}^n)\)（增广理想 \(\mathbf{I}\) 的各次幂）组成的滤过（\(\S 4\)，习题 6）。
\end{enumerate}

\begin{enumerate}
\def\labelenumi{(\alph{enumi})}
\item
  令 \(\tilde{g}\colon \mathrm{K}[\mathrm{F}] \to \tilde{\mathrm{A}}(\mathrm{X})\) 为延拓 \(g\colon \mathrm{F} \to \tilde{\mathrm{A}}(\mathrm{X})^*\) 的唯一代数同态。证明 \(\tilde{g}\) 将 \(\mathrm{I}^{\mathfrak{n}}\) 映到理想 \(\tilde{\mathrm{A}}_{\mathfrak{n}}(\mathrm{X})\) 中（参见~第 1 号），并且在取商后定义同构 \(\tilde{g}_{n}\)，其从 \(\mathrm{K}[\mathrm{F}]/\mathrm{I}^{\mathfrak{n}}\) 到 \(\tilde{\mathrm{A}}(\mathrm{X})/\tilde{\mathrm{A}}_{\mathfrak{n}}(\mathrm{X})\)。（利用 \(\tilde{g}_{n}\) 的逆同态，借助从 \(\mathrm{A}(\mathrm{X})\) 到 \(\mathrm{K}[\mathrm{F}]\) 的同态，该同态将 \(x\) 映为 \(x - 1\)，对所有 \(x \in \mathrm{X}\) 成立。）推出 \(\tilde{\mathrm{A}}(\mathrm{X})\) 同构于 \(\mathrm{K}[\mathrm{F}]\) 关于由 \((\mathrm{I}^{\mathfrak{n}})\) 定义的拓扑的豪斯多夫完备化。
\item
  假设 \(\mathbf{K} = \mathbf{Z}\)。证明滤过 \((\mathrm{I}^{\mathrm{n}})\) 是分离的（使用命题 3 和 §4 的习题 6 (e)），并且由它定义的 \(\mathbf{F}\) 的滤过（§4 的习题 6 (a)）与滤过 \((\mathbf{C}^{\mathrm{n}}\mathbf{F})\) 重合。推出 \(\tilde{g}\colon \mathbf{Z}[\mathrm{F}] \to \hat{\mathrm{A}}_{\mathbf{Z}}(\mathrm{X})\) 是单射。
\end{enumerate}

\begin{enumerate}
\def\labelenumi{\arabic{enumi}.}
\setcounter{enumi}{1}
\tightlist
\item
  设 G 是一个群，K{[}G{]} 是它在 K 上的代数，I 是其增广理想（§4，习题 6）。令 \(\varepsilon\) 表示从 K{[}G{]} 到 K 上的典范同态；于是 \(\operatorname{Ker}(\varepsilon) = I\)。
\end{enumerate}

\begin{enumerate}
\def\labelenumi{(\alph{enumi})}
\tightlist
\item
  设 M 是左 K{[}G{]}-模，Z(G, M) 是从 G 到 M 的交叉同态群（代数，第一章，§6，习题 7）。若 \(f \in \operatorname{Hom}_{\mathbb{K}[\mathbb{G}]}(\mathrm{I}, \mathrm{M})\)，令 \(f_1\) 为映射 \(g \mapsto f(g - 1)\)（从 \(\mathbf{G}\) 到 \(\mathbf{M}\)）。证明 \(f_1\) 是交叉同态（利用恒等式
\end{enumerate}

\[
g g ^ {\prime} - 1 = g (g ^ {\prime} - 1) + g - 1
\]

证明 \(f_{1}(gg') = g \cdot f_{1}(g') + f_{1}(g)\)，并证明 \(f \mapsto f_{1}\) 是从 \(\operatorname{Hom}_{\mathbb{K}[G]}(\mathbf{I}, \mathbf{M})\) 到 \(\mathbf{Z}(\mathbf{G}, \mathbf{M})\) 的同构。

\begin{enumerate}
\def\labelenumi{(\alph{enumi})}
\setcounter{enumi}{1}
\item
  取 G 为自由群 \(\mathrm{F} = \mathrm{F}(\mathrm{X})\)。证明对每个映射 \(\theta: \mathrm{X} \to \mathrm{M}\)，存在唯一的元素 \(f_{\theta}\)，属于 \(\mathrm{Z}(\mathrm{M}, \mathrm{F})\) 且延拓 \(\theta\)。（利用将 \(\mathrm{Z}(\mathrm{F}, \mathrm{M})\) 解释为 F 与 M 的半直积，参见~代数第一章上述位置。）利用 (a) 推出，族 \((x - 1)_{x \in \mathbb{X}}\) 是 I 作为左 K{[}F{]}-模的一组基。
\item
  由 (b)，每个元素 \(u \in \mathbf{K}[\mathbf{F}]\) 都能唯一写成如下形式：
\end{enumerate}

\[
u = \varepsilon (u) + \sum_ {x \in \mathbf {X}} \mathrm{D} _ {x} (u) (x - 1),
\]

其中 \(\mathrm{D}_{x}(u)\in\mathrm{K}[\mathrm{F}]\)，且除有限多个外，\(\mathrm{D}_{x}(u)\) 均为零。映射 \(u\mapsto\mathrm{D}_{x}(u)\) 称为关于 x 的偏导数。它由下列性质刻画：

\begin{enumerate}
\def\labelenumi{(\roman{enumi})}
\item
  \(D_{x}\) 是从 K{[}F{]} 到 K{[}F{]} 的 K-线性映射。
\item
  \(\mathrm{D}_{x}(uv)=u.\mathrm{D}_{x}(v)+\varepsilon(v).\mathrm{D}_{x}(u)\)，其中 u、v 属于 K{[}F{]}。
\item
  \(\mathrm{D}_{x}(x)=1\)，且 \(\mathrm{D}_{x}(y)=0\) 当 \(y\in X-\{x\}\) 时。
\end{enumerate}

若 \(n\geqslant 1\)，则

\[
\begin{array}{c} \mathrm{D} _ {x} (x ^ {n}) = 1 + x + \dots + x ^ {n - 1} \\ \mathrm{D} _ {x} (x ^ {- n}) = - x ^ {- n} - x ^ {- n + 1} - \dots - x ^ {- 1}. \end{array}
\]

\begin{enumerate}
\def\labelenumi{(\alph{enumi})}
\setcounter{enumi}{3}
\tightlist
\item
  令 \(\varepsilon_{\mathrm{A}}\) 为从 A(X) 到 K 的典范同态。证明 A(X) 的每个元素 \(u\) 都能唯一写成如下形式：
\end{enumerate}

\[
u = \varepsilon_ {\mathrm{A}} (u) + \sum_ {x \in \mathbf {X}} d _ {x} (u). x,
\]

其中 \(d_{x}(u) \in \mathrm{A}(\mathrm{X})\)，且除有限多个外，\(d_{x}(u)\) 均为零。证明 \(d_{x}: u \to d_{x}(u)\) 是 \(\mathrm{A}(\mathrm{X})\) 的自同态，可以通过连续性延拓到 \(\hat{\mathrm{A}}(\mathrm{X})\)，并具有类似于上面 (i)、(ii)、(iii) 的性质。证明 \(d_{x} \circ \tilde{g} = \tilde{g} \circ D_{x}\)，其中 \(\tilde{g}\) 是延拓 g 的从 K{[}F{]} 到 \(\hat{\mathrm{A}}(\mathrm{X})\) 的同态（参见~习题 1）。†

¶ 3. 沿用习题 2 的记号。

\begin{enumerate}
\def\labelenumi{(\alph{enumi})}
\tightlist
\item
  设 \(J\) 是 \(K[F]\) 的右理想，且包含于 \(I\)，并设 \(u\) 是 \(I\) 的一个元素。证明
\end{enumerate}

\[
u \in J. I \Leftrightarrow D _ {x} (u) \in J \quad \text { for   all } x \in X.
\]

特别地，I 的元素属于 \(\mathbf{I}^n\)（\(n \geqslant 1\)），当且仅当它的所有偏导数都属于 \(\mathbf{I}^{n-1}\)。


\begin{enumerate}
\def\labelenumi{(\alph{enumi})}
\setcounter{enumi}{1}
\tightlist
\item
  令 R 是 F 的正规子群，令 G = F/R，令 J 是典范同态 \(\gamma: K[F] \to K[G]\) 的核。证明 J 作为左理想（相应地，右理想）由元素 r - 1（其中 \(r \in R\)）生成。证明下列序列正合：
\end{enumerate}

\[
0 \longrightarrow \mathrm{J} / (\mathrm{J}. \mathrm{I}) \stackrel {\alpha} {\longrightarrow} \mathrm{I} / (\mathrm{J}. \mathrm{I}) \stackrel {\beta} {\longrightarrow} \mathrm{K} [ \mathrm{G} ] \stackrel {\delta} {\longrightarrow} \mathrm{K} \longrightarrow 0,\tag{*}
\]

其中 \(\alpha\) 由 J 包含于 I 的包含映射取商得到，\(\beta\) 由 \(\gamma\) 在 I 上的限制取商得到。

\begin{enumerate}
\def\labelenumi{(\alph{enumi})}
\setcounter{enumi}{2}
\item
  对 \(u \in \mathbf{I}\)、\(x \in \mathbf{X}\)，令 \(\overline{\mathrm{D}}_x(u)\) 表示 \(\mathrm{D}_x(u)\) 在 \(\mathrm{K}[\mathrm{G}]\) 中的像（由 \(\gamma\) 给出）。利用 (a) 证明，族 \((\overline{\mathrm{D}}_x)_{x \in \mathbf{X}}\) 在取商后定义了从 \(\mathrm{I} / (\mathrm{J}. \mathrm{I})\) 到 \(\mathrm{K}[\mathrm{G}]^{(\mathbb{X})}\) 的同构。
\item
  对 \(r \in \mathbb{R}\)，令 \(\theta(r)\) 表示 \(r - 1\) 在 \(J/(J.I)\) 中的像。证明
\end{enumerate}

\[
\theta (r r ^ {\prime}) = \theta (r) + \theta (r ^ {\prime}) \quad \text { and } \quad \theta (y r y ^ {- 1}) = y. \theta (r) \quad \text { for   } r, r ^ {\prime} \text {   in   } R, y \in F.
\]

（使用恒等式

\[
\begin{array}{c} {r r ^ {\prime} - 1 = (r - 1) (r ^ {\prime} - 1) + (r - 1) + (r ^ {\prime} - 1)} \\ {y r y ^ {- 1} - 1 = y (r - 1) (y ^ {- 1} - 1) + y (r - 1).)} \end{array}
\]

推出 \(\theta\) 定义了一个从 \(R/(R, R)\) 到 \(J/(J, I)\) 的同态，并且与 \(G\) 的作用相容；证明该同态的像生成 K-模 \(J/(J, I)\)。当 \(K = Z\) 时，证明由此得到从 \(R/(R, R)\) 到 \(J/(J, I)\) 的同构（直接定义逆同态）。（e）令 \((r_{\alpha})_{\alpha \in A}\) 是 \(R\) 中的一个族，且它生成 \(R\) 作为 \(F\) 的正规子群。证明 \(\theta(r_{\alpha})\) 生成 K{[}G{]}-模 \(J/(J, I)\)。

\begin{enumerate}
\def\labelenumi{(\alph{enumi})}
\setcounter{enumi}{5}
\tightlist
\item
  矩阵 \((\overline{\mathrm{D}}_x(r_\alpha))_{x\in \mathbf{X},\alpha \in \mathbf{A}}\) 定义了一个同态
\end{enumerate}

\[
\rho \colon \mathrm{K} [ \mathrm{G} ] ^ {(\mathrm{A})} \rightarrow \mathrm{K} [ \mathrm{G} ] ^ {(\mathrm{X})}
\]

它是左 K{[}G{]}-模的同态。证明下列序列正合：

\[
\mathrm{K} [ \mathrm{G} ] ^ {(\mathrm{A})} \xrightarrow {\rho} \mathrm{K} [ \mathrm{G} ] ^ {(\mathrm{X})} \xrightarrow {\delta} \mathrm{K} [ \mathrm{G} ] \xrightarrow {\varepsilon} \mathrm{K} \longrightarrow 0\tag{**}
\]

其中 \(\delta\) 是同态 \((u_x) \mapsto \sum_{x \in X} u_x(\gamma(x) - 1)\)。

（利用上面的 (c)、(d)、(e) 变换正合序列 (*)。）†

\begin{enumerate}
\def\labelenumi{\arabic{enumi}.}
\setcounter{enumi}{3}
\tightlist
\item
  对所有 \(n \in \mathbf{N}\)，令 \(\binom{\mathrm{T}}{n}\) 表示多项式
\end{enumerate}

\[
\mathrm{T} (\mathrm{T} - 1) \dots (\mathrm{T} - n + 1) / n!.
\]

若 \(f(\mathrm{T})\) 是多项式，令 \(\Delta f\) 表示多项式 \(f(\mathrm{T} + 1) - f(\mathrm{T})\)。

于是 \(\Delta\binom{\mathrm{T}}{0} = \Delta 1 = 0\)，并且 \(\Delta\binom{\mathrm{T}}{n} = \binom{\mathrm{T}}{n - 1}\) 在 \(n \geqslant 1\) 时成立。

\begin{enumerate}
\def\labelenumi{(\alph{enumi})}
\item
  令 \(f \in \mathbf{Q}[\mathrm{T}]\)。证明以下性质等价：
\item
  \(f\) 将 \(\mathbf{Z}\) 映入自身。
\end{enumerate}

\begin{enumerate}
\def\labelenumi{(\roman{enumi})}
\setcounter{enumi}{1}
\item
  \(f\) 是系数属于 \(\mathbf{Z}\) 的各个 \(\binom{\mathrm{T}}{n}\) 的线性组合。
\item
  \(f(0)\in \mathbf{Z}\)，且 \(\Delta f\) 将 \(\mathbf{Z}\) 映入自身。
\end{enumerate}

（对 \(\deg(f)\) 归纳，并注意 \(\deg(\Delta f) = \deg(f) - 1\) 在 \(\deg(f) \neq 0\) 时成立。）

满足上述性质的多项式 \(f\) 称为二项式多项式。两个二项式多项式的和、积和复合仍是二项式多项式。

\begin{enumerate}
\def\labelenumi{(\alph{enumi})}
\setcounter{enumi}{1}
\item
  令 \(p\) 是素数，并令 \(f\) 是二项式多项式。证明 \(f\) 将环 \(\mathbf{Z}_p\)（\(p\)-进整数环）映入自身（利用 \(f\) 的连续性以及 \(\mathbf{Z}\) 在 \(\mathbf{Z}_p\) 中稠密这一事实）。
\item
  令 \(\mathbf{P}\) 是素数集合，\(\mathrm{S}_P\) 是 P-整数的集合（§4，习题 14），令 \(\mathbf{Z}_P = \mathrm{S}_{\mathrm{P}}^{-1}\mathbf{Z}\)。证明若 \(f\) 是二项式多项式，则 \(f\) 将 \(\mathbf{Z}_{\mathrm{P}}\) 映入自身（将 (b) 应用于 \(p\) 不属于 \(\mathbf{P}\) 的素数）。
\end{enumerate}

\begin{enumerate}
\def\labelenumi{\arabic{enumi}.}
\setcounter{enumi}{4}
\tightlist
\item
  令 \(\Gamma = \Gamma(X)\) 是 \(X\) 上的 Magnus 群，令 \((\Gamma_n)\) 是其自然滤过（第 2 号）。证明与此滤过相应的分次李代数 \(\mathrm{gr}(\Gamma)\) 同构于李子代数 \(\bigoplus_{n\geqslant 1}A_n(X)\)，它是 \(A(X)\) 的子代数。若 \(X\neq \emptyset\)，此李代数不由其次数为 1 的元素生成；推出 \((\Gamma_n)\) 不是 \(\Gamma\) 的下中心列。
\end{enumerate}

¶ 6. 令 P 是素数集合，令 \(S_P\) 为 P-整数的集合（§4，习题 14），令 \(Z_P = S_P^{-1}Z\)。

\begin{enumerate}
\def\labelenumi{(\alph{enumi})}
\tightlist
\item
  假设对所有 \(s \in S_{P}\)，映射 \(k \mapsto sk\) 都是 K 到自身的双射；这等价于说 K 可以赋予 \(Z_{P}\)-代数结构。
\end{enumerate}

沿用习题 5 的记号，证明对所有 \(s \in S_{\mathbb{P}}\)，映射 \(a \mapsto a^s\) 将 \(\Gamma\) 双射到自身，并且商群 \(\Gamma / \Gamma_n\) 对每个 \(n \geqslant 1\) 也具有同样性质。若 \(t \in \mathbf{Z}_{\mathbb{P}}\) 且 \(a \in \Gamma\)（相应地，\(a \in \Gamma / \Gamma_n\)），则按 §4 习题 16 定义 \(a^s\)；将 \(a\) 写成 \(1 + \alpha\) 的形式，其中 \(\alpha \in \hat{\mathrm{A}}_1(\mathrm{X})\)，证明

\[
a ^ {t} = (1 + \alpha) ^ {t} = \sum_ {n = 0} ^ {\infty} \binom {t} {n} \alpha^ {n}.
\]

（注意系数 \(\binom{t}{n}\) 属于 \(\mathbf{Z}_{\mathrm{P}}\)，参见~习题 4。）

\begin{enumerate}
\def\labelenumi{(\alph{enumi})}
\setcounter{enumi}{1}
\item
  现在假设 \(\mathbf{K} = \mathbf{Z}_{\mathrm{P}}\)。令 \(c\) 是满足 \(\geqslant 1\) 的整数。将群 \(\mathbf{F}^c = \mathbf{F} / \mathbf{C}^c\mathbf{F}\) 与 \(\Gamma /\Gamma_c\) 的一个子群相识别，所用的是由 \(g\) 取商得到的同态（定理 2 的记号）。群 \(\Gamma / \Gamma_c\) 是类 \(< c\) 的 P-无挠、P-可除群。令 \(F_P^c\) 为 \(F^c\) 在 \(\Gamma / \Gamma_c\) 中的 P-饱和（§4，习题 14），并令 \(i\) 是从 \(F^c\) 到 \(F_P^c\) 的嵌入。序对 \((i, F_P^c)\) 是 \(F^c\) 的 P-包络（§4，习题 15）。推出每个幂零群都有 P-包络（注意每个类 \(< c\) 的幂零群都是某个适当 X 的群 \(F^c\) 的商，并使用§4 习题 15 的 (b) 部分）。
\item
  若 \(n \leqslant c\)，令 \(\mathbf{F}_{\mathbf{P},n}^{\mathrm{c}}\) 为 \(\mathbf{F}_{\mathbf{P}}^{\mathrm{c}}\) 与 \(\Gamma_{n}/\Gamma_{c}\) 的交；若 \(n \geqslant c\)，记 \(\mathbf{F}_{\mathbf{P},n}^{\mathrm{c}} = \{e\}\)。证明 \(\mathbf{F}_{\mathbf{P},n}^{\mathrm{c}}\) 是 \(\mathbf{C}^{\mathrm{c}}\mathbf{F}^{\mathrm{c}}\) 在 \(\Gamma/\Gamma_{c}\) 中的 P-饱和。滤过 \((\mathbf{F}_{\mathbf{P},n}^{\mathrm{c}})\) 是 \(\mathbf{F}_{\mathbf{P},n}^{\mathrm{c}}\) 的整中心滤过；令 \(\text{gr}(\mathbf{F}_{\mathbf{P}}^{\mathrm{c}})\) 为与此滤过相伴的分次模。证明若 \(n < c\)，则 \(\text{gr}_{\mathbf{n}}(\mathbf{F}_{\mathbf{P}}^{\mathrm{c}})\) 在 \(\text{gr}_{\mathbf{n}}(\Gamma) = \mathbf{A}_{\mathbf{Z}_{\mathbf{P}}}^{n}(\mathbf{X})\) 中的像是 \(\mathbf{S}_{\mathbf{P}}^{-1}. \mathbf{L}_{\mathbf{Z}}^{n}(\mathbf{X}) = \mathbf{L}_{\mathbf{Z}_{\mathbf{P}}}^{n}(\mathbf{X})\)。推出李代数 \(\text{gr}(\mathbf{F}_{\mathbf{P},n}^{\mathrm{c}})\) 由其次数为 1 的元素生成，从而 \(\mathbf{F}_{\mathbf{P},n} = \mathbf{C}^{n}(\mathbf{F}_{\mathbf{P}}^{\mathrm{c}})\) 对所有 \(n\) 都成立（\(\S 4\)，习题 1）。证明群 \(\mathbf{F}_{\mathbf{P}}^{\mathrm{c}}\) 由 \(x^{1/s}\)（其中 \(x \in \mathbf{X}\)、\(s \in \mathbb{S}_{\mathbf{P}}\)）生成（注意这些元素在 \(\text{gr}_1(\mathbf{F}_{\mathbf{P}}^{\mathrm{c}}) \simeq \mathbf{Z}_{\mathbf{P}}^{(X)}\) 中的像生成群 \(\text{gr}_1(\mathbf{F}_{\mathbf{P}}^{\mathrm{c}})\)，并应用 \(Algebra\) 第一章 §6 第 3 号命题 8 的推论 3）。
\item
  令 \(\mathrm{H}\) 是相对于 \(\mathbf{X}\) 的 Hall 集（§2，第 10 号），令 \(\mathrm{H}(c)\) 为由 \(\mathrm{H}\) 中长度 \(< c\) 的元素组成的子集。对 \(m \in \mathrm{H}(c)\)，令 \(\phi_c(m)\) 表示基本换位子在 \(\mathbf{F}^c\) 中的像；这里的基本换位子是 \(\phi(m)\)，由 \(m\) 定义（参见~第 4 号的注）。证明对所有 \(w \in \mathbf{F}_P^c\)，存在唯一的 \(\alpha \in \mathbf{Z}_{\mathbf{P}}^{\mathbf{H}(c)}\)，使得
\end{enumerate}

\[
w = \prod_ {m \in H (c)} \phi_ {c} (m) ^ {\alpha (m)}.
\]

（使用上面得到的 gr(F \(_{P}^{c}\) ) 的刻画。）

\begin{enumerate}
\def\labelenumi{\arabic{enumi}.}
\setcounter{enumi}{6}
\tightlist
\item
  沿用习题 6 的记号。
\end{enumerate}

\begin{enumerate}
\def\labelenumi{(\alph{enumi})}
\item
  令 G 是类 \(<c\) 的幂零群，令 \((i, \overline{G})\) 是 G 的一个 P-包络。令 \(\rho: \mathrm{F}^{c} \to \mathrm{G}\) 是满同态（适当选取 X 时这样的同态存在），令 \(\bar{\rho}\) 是相应的 \(\mathrm{F}_{\mathbf{P}}^{c} \in \overline{G}\) 的同态（§4，习题 15）。同态 \(\bar{\rho}\) 是满射，并将 \(\mathrm{C}^{n}\mathrm{F}_{\mathbf{P}}^{c} = \mathrm{F}_{\mathbf{P}, n}^{c}\) 映到 \(\mathrm{C}^{n}\overline{G}\) 上。证明 \(\mathrm{C}^{n}\overline{G}\) 是 \(i(\mathrm{C}^{n}\mathrm{G})\) 在 \(\overline{G}\) 中的 P-包络，并且李代数 \(\mathrm{gr}(\overline{G})\) 可与 \(\mathbf{Z}_{\mathbf{P}} \otimes \mathrm{gr}(\mathrm{G})\) 相识别。
\item
  推出 G 是 P-可除且 P-无挠的，当且仅当各个 \(C^{n}G/C^{n+1}G\) 具有这些性质。
\end{enumerate}

¶ 8. 设 K = Z 且 X 有限。对每个整数 \(k \geqslant 0\)，令 \((e_{k,\alpha})\) 表示 \(A_k(X)\) 的一组基。

\begin{enumerate}
\def\labelenumi{(\alph{enumi})}
\tightlist
\item
  令 \((w_{n})_{n\in \mathbf{N}}\) 是 F 中元素的一个序列。令 \(w_{k,\alpha}(n)\) 为 \(e_{k,\alpha}\) 的系数；该系数取自 \(g(w_{n})\in \hat{\mathrm{A}} (\mathrm{X})\) 中次数为 \(k\) 的项。则：
\end{enumerate}

\[
g (w _ {n}) = \sum_ {k, \alpha} w _ {k, \alpha} (n) e _ {k, \alpha} \quad \text { for   all } n \in \mathbf {Z}.
\]

序列 \((w_{n})\) 称为典型的，如果对每个有序对 \((k,\alpha)\)，函数 \(w_{k,\alpha}\colon n\mapsto w_{k,\alpha}(n)\) 是次数 \(\leqslant k\) 的二项式多项式，参见~习题 4。此条件与基 \((e_{k,\alpha})\) 的选择无关。证明它等价于存在一个序列 \((a_{k})_{k\in \mathbf{N}}\)，其元素属于 \(\hat{\mathbf{A}}_{\mathbf{q}}(\mathbf{X})\)，使得 \(\omega (a_k)\geqslant k\) 对所有 \(k\) 都成立，并且

\[
g (w _ {n}) = \sum_ {k = 0} ^ {\infty} n ^ {k} a _ {k} \quad \text { for   all } n \in \mathbf {Z}.
\]

\begin{enumerate}
\def\labelenumi{(\alph{enumi})}
\setcounter{enumi}{1}
\item
  证明若 \((w_{n})\) 和 \((w_{n}^{\prime})\) 是两个典型序列，则 \((w_{n}^{-1})\) 和 \((w_{n}w_{n}^{\prime})\) 也都是典型序列。
\item
  令 \(w\) 是 \(\mathbf{C}^k\mathbf{F}\) 中的元素，其中 \(k \geqslant 1\)，令 \(f\) 是次数 \(\leqslant k\) 的二项式多项式。证明 \((w^{f(n)})\) 是典型序列。特别地，若 \(z \in \mathbf{F}\)，则幂序列 \((z^n)\) 是 \(z\) 的典型幂序列。
\item
  令 \((w_{n})\) 是 F 中元素的一个序列。证明 \((w_{n})\) 是典型的，当且仅当存在 F 中的 \(y_{0}, y_{1}, \ldots, y_{n}, \ldots\)，满足 \(y_{i} \in \mathbf{C}^{i}\mathbf{F}\) 对所有 \(i \geqslant 1\) 都成立，并且
\end{enumerate}

\[
w _ {n} \equiv y _ {0} y _ {1} ^ {n} \dots y _ {k} ^ {\binom {n} {k}} \mod \mathrm{C} ^ {k + 1} \mathrm{F}
\]

对所有 \(n \in \mathbf{Z}\) 和所有 \(k \geqslant 1\) 均成立。（确定 \(y_{i}\)：由 \((w_{n})\) 依次令 \(n = 0, 1, \ldots\)。例如 \(w_{0} = y_{0}, w_{1} = y_{0}y_{1}, \ldots\)。）

\begin{enumerate}
\def\labelenumi{(\alph{enumi})}
\setcounter{enumi}{4}
\tightlist
\item
  令 H 是相对于 X 的 Hall 集，\((w_{n})\) 是 F 中元素的一个序列。令 \((\alpha(m,n))_{m\in\mathbb{H},n\in\mathbf{Z}}\) 为满足下式的整数族：
\end{enumerate}

\[
w _ {n} \equiv \prod_ {m \in \mathrm{H}} \phi (m) ^ {\alpha (m, n)} \mod \mathrm{C} ^ {c} \mathrm{F}
\]

对所有 \(n \in \mathbf{Z}\) 和所有 \(c \geqslant 1\) 均成立（参见~第 4 号的注）。证明 \((w_n)\) 是典型的，当且仅当对所有 \(m \in \mathbf{H}\)，函数 \(n \mapsto \alpha(m, n)\) 是次数 \(\leqslant l(m)\) 的二项式多项式，其中 \(l(m)\) 是 \(m\) 的长度。†

\begin{enumerate}
\def\labelenumi{(\alph{enumi})}
\setcounter{enumi}{5}
\tightlist
\item
  序列 \((w_{n})\) 称为 1-典型，如果相应的函数 \(w_{k,\alpha}\) 是次数 \(\leqslant k - 1\) 的二项式多项式。证明若 \((w_{n})\) 是 1-典型的，则存在 \(y_{i} \in \mathbf{C}^{i + 1}\mathbf{F}\)（\(i = 0, 1, \ldots\)），使得
\end{enumerate}

\[
w _ {n} \equiv y _ {0} y _ {1} ^ {n} \dots y _ {k} ^ {\binom {n} {k}} \mod \mathrm{C} ^ {k + 1} \mathrm{F}
\]

对所有 \(n \in Z\) 和所有 \(k \geqslant 1\) 均成立。

¶ 9. 令 X 是含两个元素 \{x, y\} 的集合。

\begin{enumerate}
\def\labelenumi{(\alph{enumi})}
\tightlist
\item
  证明存在一个序列 \(w_{2}, w_{3}, \ldots\)，其元素属于 \(\mathbf{F}\)，且 \(w_{i} \in \mathbb{C}^{i}\mathbb{F}\) 对所有 \(i\) 都成立，并且
\end{enumerate}

\[
(x y) ^ {n} \equiv x ^ {n} y ^ {n} w _ {2} ^ {\binom {n} {2}} \dots w _ {i} ^ {\binom {n} {i}} \mod \mathrm{C} ^ {i + 1} \mathrm{F}
\]

对所有 \(n \in \mathbf{Z}\) 和所有 \(i \geqslant 1\) 均成立。（将习题 8（\(d\)）应用于序列 \(x^{-n}(xy)^n\)。）证明 \(w_2 = y^{-1}(y^{-1}, x^{-1})y \equiv (x, y)^{-1} \mod \mathrm{C}^3\mathrm{F}\)。

\begin{enumerate}
\def\labelenumi{(\alph{enumi})}
\setcounter{enumi}{1}
\tightlist
\item
  设 G 是类 \(\leqslant c\) 的幂零群，且 x、y 属于 G。由上文推出下列公式（“Hall 公式”）：
\end{enumerate}

\[
(x y) ^ {n} = x ^ {n} y ^ {n} \prod_ {i = 2} ^ {c} w _ {i} (x, y) ^ {\binom {n} {i}}.
\]

\begin{enumerate}
\def\labelenumi{(\alph{enumi})}
\setcounter{enumi}{2}
\tightlist
\item
  令 \(p\) 为素数。证明存在 \(u_{i} \in \mathbf{C}^{i}\mathbf{F}\)（ \(2 \leqslant i \leqslant p - 1\) ）以及 \(w \in \mathbf{C}^{p}\mathbf{F}\)，使得
\end{enumerate}

\[
(x y) ^ {p} = x ^ {p} y ^ {p} u _ {2} ^ {p} \dots u _ {p - 1} ^ {p} w.
\]

（使用 (a)，注意 \(\binom{p}{i}\) 被 \(p\) 整除，只要 \(1 < i < p\)。）令 \(\bar{w}\) 为 \(w\) 在 \(\mathrm{gr}_p(\mathsf{F}) = \mathrm{L}_{\mathbf{Z}}^p (\mathbf{X})\) 中的像，令 \(\bar{w}_p\) 为 \(\bar{w}\) 在 \(\mathrm{L}_{\mathbf{F}_p}^p (\mathbf{X})\) 中的像。证明 \(\bar{w}_p = \Lambda_p(x,y)\)，参见~第一章，§ 1，习题 19。（将 F 通过同态 \(g\) 延拓到 \(\tilde{\mathbf{A}}(\mathbf{X})\)，并比较次数为 \(p\) 的 \(g((xy)^p)\) 与 \(g(x^p y^p u_2^p \ldots u_{p-1}^p w)\) 中的项。）前者等于 \((x + y)^p\)，后者模 \(p\) 同余于 \(x^p + y^p + \bar{w}_p\)。故结论成立。）

证明存在 \(a \in \mathbf{C}^2\mathbf{F}\) 以及 \(w' \in \mathbf{C}^p\mathbf{F}\)，使得

\[
(x y a) ^ {p} = x ^ {p} y ^ {p} w ^ {\prime}.
\]

（使用上式计算 \((xy)^p\)。）由此推出，类为 \(p\) 的幂零群中，\(p\) 次幂构成一个子群。

\begin{enumerate}
\def\labelenumi{(\alph{enumi})}
\setcounter{enumi}{3}
\tightlist
\item
  证明存在一个序列 \(v_{3}, v_{4}, \ldots\)，其元素属于 \(\mathbf{F}\)，且 \(v_{i} \in \mathbf{C}^{i}\mathbf{F}\) 对所有 \(i\) 都成立，并且
\end{enumerate}

\[
(x ^ {n}, y) \equiv (x, y) ^ {n} v _ {3} ^ {\binom {n} {2}} \dots v _ {i + 1} ^ {\binom {n} {i}} \mod \mathrm{C} ^ {i + 2} \mathrm{F}
\]

对所有 \(n \in \mathbf{N}\) 和所有 \(i \geqslant 1\) 均成立。（将习题 8(f) 应用于 \((x^n, y)\) 的序列。）

推出，若 \(p\) 为素数，则存在 \(t_i \in \mathrm{C}^i\mathrm{F}\)（ \(3 \leqslant i \leqslant p\)）以及 \(z \in \mathrm{C}^{p+1}\mathrm{F}\)，使得

\[
(x ^ {p}, y) = (x, y) ^ {p} t _ {3} ^ {p} \dots t _ {p} ^ {p} z.
\]

证明 \(\bar{z}_p\) 是 \(z\) 在 \(\mathrm{L}_{\mathbf{F}_p}^{p + 1}(\mathbf{X})\) 中的像，且等于 (ad \(x)^p (y)\)。（方法同 \((c).)\)

¶ 10. 令 p 为素数，令 G 是带有实中心滤过 (G \(_{\alpha}\) ) 的群。若关系 x ∈ G \(_{\alpha}\) 蕴含 x \(^{p}\) ∈ G \(_{p\alpha}\)，则称滤过 (G \(_{\alpha}\) ) 为受限滤过。

\begin{enumerate}
\def\labelenumi{(\alph{enumi})}
\item
  证明李代数 \(\operatorname{gr}(\mathbf{G})\)（与 \((\mathbf{G}_{\alpha})\) 相关）满足 \(p.\operatorname{gr}(\mathbf{G}) = 0\)，因而可以赋予 \(\mathbf{F}_p\) 上的代数结构。
\item
  令 \(\xi \in \mathrm{gr}_{\alpha}(\mathbf{G})\)，令 \(x\) 是 \(\xi\) 在 \(\mathbf{G}_{\alpha}\) 中的一个代表元。证明 \(x^{p}\) 在 \(\mathrm{gr}_{p\alpha}(\mathbf{G})\) 中的像与 \(x\) 的选取无关（使用习题 9(c)）。若将此像记为 \(\xi^{[p]}\)，证明 \(\xi \mapsto \xi^{[p]}\) 是从 \(\mathrm{gr}_{\alpha}(\mathbf{G})\) 到 \(\mathrm{gr}_{p\alpha}(\mathbf{G}')\) 的线性映射，并且 \((\xi + \xi')^{[p]} = \xi^{[p]} + \xi'^{[p]} + \Delta_p(\xi, \xi')\)。（方法相同。）
\end{enumerate}

证明，若 \(\xi \in \mathrm{gr}_{\alpha}(\mathbf{G})\) 且 \(\eta \in \mathrm{gr}_{\beta}(\mathbf{G})\)，则

\[
[ \xi^ {[ p ]}, \eta ] = (\operatorname{ad} \xi) ^ {p} (\eta).
\]

（使用习题 9(d)。）

\begin{enumerate}
\def\labelenumi{(\alph{enumi})}
\setcounter{enumi}{2}
\tightlist
\item
  证明存在唯一的 \(p\)-映射，它是从 \(\operatorname{gr}(\mathbf{G})\) 到自身的映射（第一章，§ 1，习题 20），且延拓了映射 \(\xi \mapsto \xi^{[p]}\)，该映射定义在各个 \(\operatorname{gr}_{\alpha}(\mathbf{G})\) 上。赋予这一结构后，\(\operatorname{gr}(\mathbf{G})\) 是一个李 \(p\)-代数，称为与受限滤过 \((\mathbf{G}_{\alpha})\) 相关的李代数。
\end{enumerate}

¶ 11. (a) 令 A 是满足 §4 第 5 号条件的滤过代数，令 \(\Gamma = A^{*} \cap (1 + A_{0}^{+})\)。在 A 上的滤过诱导出 \(\Gamma\) 上的滤过（同上，命题 2）。证明若 K 是特征为 \(p > 0\) 的域，则 \((\Gamma_{\alpha})\) 是受限滤过（习题 10），并且同上命题 3 所定义的 \(\mathrm{gr}(\Gamma)\) 到 \(\mathrm{gr(A)}\) 的嵌入与李 \(p\)-代数结构相容；这些结构分别定义在 \(\mathrm{gr}(\Gamma)\) 和 \(\mathrm{gr(A)}\) 上。

\begin{enumerate}
\def\labelenumi{(\alph{enumi})}
\setcounter{enumi}{1}
\item
  假设 \(\mathbf{K} = \mathbf{F}_p\)。将上述结果应用于滤过代数 \(\hat{\mathbf{A}}(\mathbf{K})\)，其中 \(\mathrm{F} = \mathrm{F}(\mathrm{X})\) 借助同态 \(g\) 嵌入。滤过 \((\Gamma_n)\) 的 \(\Gamma\) 诱导出滤过 \((\mathbb{F}_n)\)，该滤过定义在 \(\mathbf{F}\) 上，这是一个受限滤过。李 \(p\)-代数 \(\operatorname{gr}(\mathbf{F})\) 可识别为一个李 \(p\)-子代数 \(\mathrm{A}(\mathbf{X})\)，其中含有 \(\mathbf{X}\)。推出 \(\operatorname{gr}(\mathbf{F})\) 含有自由李 \(p\)-代数 \(\mathrm{L}(\mathbf{X}, p)\)，参见~§ 3，习题 4(e)。
\item
  令 \(\mathbf{H}\) 是关于 \(\mathbf{X}\) 的 Hall 集；对每个整数 \(i\)，令 \(\mathbf{H}_i\) 为 \(\mathbf{H}\) 中长度为 \(i\) 的元素所成的集合。令 \(w \in \mathbf{F}\)，并令 \(\alpha_i \in \mathbf{Z}^{(\mathbf{H}_i)}\) 满足
\end{enumerate}

\[
w \equiv \prod_ {i = 1} ^ {n} \prod_ {m \in H_i} \phi (m) ^ {\alpha_ {i} (m)} \mod C ^ {n + 1} F
\]

对所有 \(n\) 均成立，参见~第 4 号。对每个 \(m \in \mathbf{H}\)，令 \(l(m)\) 表示 \(m\) 的长度，令 \(q(m, w)\) 表示 \(p\) 的最大次幂，使其整除 \(\alpha_{1}(m)\)，其中 \(i = l(w)\)（若 \(\alpha_{1}(m) = 0\)，约定 \(q(m, w) = +\infty\)）。令 \(\mathrm{N} = \inf_{m \in \mathbb{H}} (l(m).q(m, w))\)。证明 \(\phi(m)^{\alpha_{1}(m)}\) 属于 \(\mathrm{F}_{\mathrm{N}}\)，并且甚至属于 \(\mathrm{F}_{\mathrm{N+1}}\)（当 \(l(w).q(m, w) > \mathrm{N}\) 时）。证明若 \(l(w).q(m, w) = \mathrm{N}\)，则 \(\phi(m)^{\alpha_{1}(m)}\) 在 \(\mathrm{gr}_{\mathrm{N}}(\hat{\mathrm{A}}(\mathbf{X})) = \mathrm{A}^{\mathrm{N}}(\mathbf{X})\) 中的像是 \(\overline{m}^{q(m, w)}\) 的非零倍数，其中 \(\overline{m}\) 是 \(\mathrm{L}(\mathbf{X})\) 中由 \(m\) 定义的元素（\(\S 2\)，第 11 号）。现在这些元素线性无关（\(\S 3\)，习题 4）；推出 \(w\) 在 \(\mathrm{gr}_{\mathrm{N}}(\mathbf{F})\) 中的像是 \(\neq 0\) 且属于 \(\mathrm{L}(\mathbf{X}, p)\)，从而 \(\mathrm{gr}(\mathbf{F}) = \mathrm{L}(\mathbf{X}, p)\)。

\begin{enumerate}
\def\labelenumi{\arabic{enumi}.}
\setcounter{enumi}{11}
\tightlist
\item
  令 g 是特征为 p \textgreater{} 0 的域 k 上的李代数。
\end{enumerate}

\begin{enumerate}
\def\labelenumi{(\alph{enumi})}
\tightlist
\item
  令 \(h\) 为满足 \(\leqslant p - 1\) 的整数。\(\mathfrak{g}\) 称为满足 Engel 的第 \(h\) 条件，是指 \((\operatorname{ad} x)^h(y) = 0\) 对每个有序对 \((x, y)\)（其元素属于 \(\mathfrak{g}\)）都成立。证明这等价于
\end{enumerate}

\[
\sum_ {\sigma \in \mathfrak {S} _ {h}} \operatorname{ad} (x _ {\sigma (1)}) \circ \dots \circ \operatorname{ad} (x _ {\sigma (h)}) = 0
\]

当 \(x_1, \ldots, x_h\) 属于 \(\mathfrak{g}\) 时成立（将公式 \((\operatorname{ad} x)^h = 0\) 应用于 \(x_i\) 的线性组合）。


\begin{enumerate}
\def\labelenumi{(\alph{enumi})}
\setcounter{enumi}{1}
\item
  令 \(x, y\) 为 \(\mathfrak{g}\) 中的元素，满足 \(\Lambda_p(ax, by) = 0\)，其中 \(a, b\) 属于 \(k\)。证明若 \(\Lambda_p^{r,s}\) 是 \(\Lambda_p\) 的双齐次分量，且双次数为 \((r, s)\)，则 \(\Lambda_p^{r,s}(x, y) = 0\) 对每个有序对 \((r, s)\) 都成立。推出（取 \(r = p - 1\)、\(s = 1\)）\((\text{ad } x)^{p-1}(y) = 0\)。特别地，若 \(\Lambda_p(x, y) = 0\)，其中 \(x, y\) 属于 \(\mathfrak{g}\)，则李代数 \(\mathfrak{g}\) 满足 Engel 的第 \((p - 1)\) 条件。
\item
  令 \(\mathfrak{c}\) 为 \(\mathfrak{g}\) 的中心。假设 \((\operatorname{ad} x)^p = 0\) 对所有 \(x \in \mathfrak{g}\) 都成立。证明 \(\mathfrak{g} / \mathfrak{c}\) 满足 Engel 的第 \((p - 1)\) 条件。（证明 \(\Lambda_p(\operatorname{ad} x, \operatorname{ad} y) = 0\) 对 \(x, y\) 属于 \(\mathfrak{g}\) 成立，并将 \((b)\) 应用于李代数 \(\operatorname{ad} \mathfrak{g} = \mathfrak{g} / c\)。）
\item
  令 G 是满足 \(x^{p}=1\) 对所有 \(x\in G\) 成立的群，令 \((\mathrm{G}_{\alpha})\) 是 G 的实中心滤过。滤过 \((\mathrm{G}_{\alpha})\) 是受限的（习题 10），而 \(\mathrm{gr}(\mathrm{G})\) 是 p-映射为零的李 p-代数。推出 \(\mathrm{gr}(\mathrm{G})\) 满足 Engel 的第 \((p-1)\) 条件。†
\end{enumerate}

